\documentclass[10pt,aspectratio=169]{beamer}
\usetheme{Warsaw}
\usepackage{amsmath,amssymb}

\title{Algebra And Number Theory Talk}
\author{Speaker}
\date{\today}

\begin{document}
\begin{frame}\titlepage\end{frame}

\begin{frame}{Objects}
  \begin{itemize}
    \item Rings, fields, groups, modules, varieties, or representations.
    \item Morphisms and equivalence relation.
    \item Characteristic, finiteness, or local/global assumptions.
  \end{itemize}
\end{frame}

\begin{frame}{Running Example}
  \[
    \mathbb{Z}/n\mathbb{Z}, \quad \mathbb{F}_q[x], \quad G \curvearrowright X.
  \]
  Explain the example before stating the general theorem.
\end{frame}

\begin{frame}{Structure Theorem}
  \begin{theorem}
    State the algebraic or arithmetic classification result.
  \end{theorem}
  \begin{itemize}
    \item Source: theorem number, proposition, or standard reference.
    \item Status: proved, conditional, conjectural, or example-driven.
    \item Ledger: record characteristic, finiteness, local/global, and category assumptions.
  \end{itemize}
\end{frame}

\begin{frame}{Proof Route}
  \begin{enumerate}
    \item Reduce to local or simple objects.
    \item Apply key lemma or descent argument.
    \item Reassemble and check uniqueness.
  \end{enumerate}
\end{frame}
\end{document}
